\newpage
\thispagestyle{empty}
\mbox{}
\newpage

%%%%%%%%%%%%%%%%%%%%%%%%%%%%%%%%%%%%%%%%%%%%%%%%%%%%%%%%%%%%

\chapter*{Remerciements}
\label{chap:remerciements}
{\fontsize{11}{13}\selectfont
The end of a thesis marks, for the PhD candidate, the end of an era.
A period filled with encounters, new knowledge, emotional highs and lows, moments of doubt and of clarity, intense work and well-deserved breaks, new friendships, inspiring discussions, and so much more.
In the following lines, I would like to acknowledge the people who have made these three years a truly special time in my life.
\vspace{\baselineskip}

Although it was\----most of the time\----a great period, I would like to begin with the four persons who have accepted to put an end to it: dear members of my jury, thank you very much for accepting to review my work, it is a great honour to be read by you.
\vspace{\baselineskip}

%Louis et JC
Ensuite, j'aimerais continuer par ceux qui l'ont faite démarrer, Louis et Jean-Christophe, mes deux directeurs de thèse.
Merci pour la confiance que vous m'avez accordée à mon arrivée au laboratoire en M2, et que vous avez sans cesse renouvelée tout au long de la thèse.
 Jean-Christophe, merci pour ton conseil scientifique spontané et efficace, ainsi que pour ta rigueur et ton sérieux dans lesquels j'ai su me retrouver. %Travailler à tes côtés .....

Un merci tout particulier à Louis : ton soutien technique et scientifique qui n'aura jamais failli, et ce, jusqu'au tout derniers moments, compte parmi les éléments les plus formateurs de mon parcours d'études.
 Ton intuition scientifique et tes idées comptent pour beaucoup dans ce qu'est devenu cette thèse.
 Tu m'a conduit sur tes pas, m'apprenant chaque jour (ou du moins à chaque compilation de modèle) ce qui m'aura permis demain de me lancer vers la suite de mon parcours de modélisateur.
Je garde chaudement en mémoire les traversées pédestres qui marquaient chacun de nos déplacements : Grenoble, Viennes, Clermont, Cambridge... À quand la traversée de Paris ?
\vspace{\baselineskip}

%Jennie et Rémy
Je souhaite remercier Jennie et Rémy, qui malgré la distance ont en quelque sorte parrainé ma thèse.
% Jennie pour la confiance que tu m'accordes ainsi que tout ce que tu m'as perms de faire.
Jennie, la confiance que tu m'accordes est un des moteurs qui me donnent la force d'avancer sur le chemin de la recherche.
Certes, peut-être que tu me faisais un peu peur au début... mais j'ai toujours su que je pouvais compter sur toi.
Plus d'une fois nos discussions en tête-à-tête auront suffit à me remettre d'aplomb, ce qui fait que c'est en partie grâce à toi que j'en suis arrivé là.

% Rémy un vrai modèle
Rémy, ton enthousiasme et ta sympathie sont les premières choses auxquelles on a affaire lorsqu'on te rencontre, mais ça ne nous prépare pas à ce que l'on gagne lorsqu'on prend ta roue. 
Alors oui, je sais, ``si tu peux le faire, tout le monde peut le faire'', n'empêche que sans la personne pour te passer le livre, difficile de trouver la réponse à la vie, l'Univers et tout le reste.
\vspace{\baselineskip}

%Les amis du bureau 307 et la génération de thésards : Nicolas Streel, Thomas L., Sébastien V., Natalie
J'aimerais remercier mes camarades du bureau 307, où nos aventures ont commencé par une arrivée de météorite qui pesait son poids, et qui ce sont prolongées tout au long de ces trois années. Merci Nicolas, Thomas et Yaël pour votre si sympathique compagnie.

Merci également aux autres doctorants et doctorantes de ma génération de thésards (et un peu au dessus), Sébastien, Natalie, Florent, Nadir, Tanguy et Clara. Vous avez constitué le premier paysage de ma thèse au LATMOS.
\vspace{\baselineskip}

%Les doctorants avec qui j'ai travaillé Thomas, Yaël, Antoine E.
Un merci particulier aux doctorants avec qui j'ai eu le plaisir de travailler. 
Thomas, malgré notre petite différence d'âge, on aura été jumeau jusqu'au bout. Bien que tu sois musicien, je sais que tu ne raffoles pas des trompettes à rubans, alors je m'en tiens là : merci pour tout ce que l'on a partagé, c'était un plaisir de faire cette thèse à tes côtés.

Antoine E., quatre heures de train filent comme du \texit{matériel roulant grande vitesse} lorsqu'on est à tes côtés. .......

Yaël,  % petit frère de thèse, les aventures ne font que démarrer, tu m'as fait confiance et maintenant je sais que je peux me reposer sur toi. Vivement la suite des aventures.

\vspace{\baselineskip}

%Les copains proches de 2025 : Nicolas T., Paul L., Maxence L., Valentine J., Flore, Louis Milhamont, Nejla E\`co, Selviga, Joris (Nadir)
Sur sa dernière année, ma thèse a pris une saveur charmante au petit goût de magnésie que je dois à Nicolas T., Paul, Valentine, Maxence, Selviga, Louis Milhamont (oui, je te mets dans le groupe des grimpeurs), Nejla, Joris, Yaël et Flore.
%Je me sens bien avec vous
% Maxence pour la natation, la vie du labo, et l'organisation des blagues
Avant de m'user les doigts sur les murs, Maxence, tu as joué le coup de maître de me faire replonger dans une piscine. Mais c'est la collection de nos innombrables \textit{farces de couloir} qui restera gravée.
\vspace{\baselineskip}

Il me tient à c\oe ur de remercier toutes les personnes de Climactions, et particulièrement Stéphanie B., Xavier C., Juliette M., Olivier A., Julie D., Marie-Alice F., Marc D. et Amélie C., pour votre sensibilité, votre écoute et votre engagement, non pas intarissable, mais si sincère. C'est à vos côtés que j'ai fait de l'IPSL ma seconde maison.
%
%Les inmanquables du LATMOS : Jacques L., Richard W., François R., Pascal, Cristelle, Gaëlle, etc.
%
%Les gens de Grenoble et d'ERCA.
%
%Les personnes de l'UFR TEB : Olivier K., Matthieu Chassé, Valérie Plagnes, Guillaume Gastineau, le troisième laron, Aymeric Spiga, et Isabelle des clefs
Je remercie les personnels de l'UFR TEB auprès de qui j'ai travaillé lors de ma mission d'enseignement : Olivier Kachnic, Valérie Plagnes, Guillaume Gastineau, François Baudin, Aymeric Spiga et Isabelle Bauwe.
Je remercie particulièrement Mathieu Chassé avec qui j'ai eu le plaisir de travailler à la préparation des supports d'exposés de L1 TCE. Ton aide, ta rigueur et ton expertise furent inspirantes, et j'espère que nous serons amenés à retravailler ensemble.
\vspace{\baselineskip}
%
%Le comité de thèse.
%
%Paul Z. and Julia K. and Annicka E.
I would like to thank Paul Z. and Julia K. for our insightful discussions about PBAPs and beyond. I look forward to working with you in the future.
\vspace{\baselineskip}
 
%Lea R
Lea, merci pour ton soutient et ton écoute, tu savais reconnaître quand le poids de la thèse pesait de travers sur mes épaules, et il suffisait d'un thé et de quelques mots pour que tu me soulages de son portage.
Parmi nos joies et nos peines, nos souffrances et nos tendresses, il y a des choses irréparables, mais je ne voudrais pas d'un monde où elles n'auraient pas existé. 
\vspace{\baselineskip}


%Niels Dutrievoz
%

%Autres ami·es du LMD : Antoine B., Enora M., Léa B., Siméon 
Je remercie également mes ami·es du LMD, et particulièrement Antoine B. et Enora M. parce que lorsqu'on est avec vous, le monde en devient plus beau. % pour les formidables personnes que vous êtes.
Et bien sûr Léa B., pour ton énergie à revendre et ton engagement qui ne sont pas près d'être oubliés en 307.
\vspace{\baselineskip}

Enfin, je remercie du fond du c\oe r mes proches.
%Chloé et Etienne
Chloé et Etienne, vous m'avez accompagné par monts et par vauts sans jamais faillir.....
% ...vous doit pour beaucoup

Merci à mon oncle Phil, pour tes références dont certaines se cachent dans ces lignes, et à Coc' pour ton soutient inconditionnel.

%Papa et Maman
Maintenant
%
%
}
